
\chapter{Warstwa programowa stacji bazowej}
\label{chap:oprogramowanie}

\section{Środowisko uruchomieniowe}
\label{sec:srodowisko}
%Ubuntu, Docker (obraz, kontener, compose), usługi systemd,
% dump1090 na hoście (porty 30002 i 30003).

% ---------------------------------------------------------------------
\section{Węzły akwizycji danych}
\label{sec:akwizycja}
% sdr_node (dwa połączenia TCP, buforowanie linii,
% automatyczne ponowne łączenie), gps_node, power_monitor_node.
\begin{figure}[!htb]
\centering
\begin{tikzpicture}[
    font=\scriptsize,
    >={Stealth[length=2mm]},
    term/.style={draw, thick, rounded corners=3.5mm, fill=akcent!25,
                 minimum height=7mm, minimum width=22mm, align=center},
    we/.style={draw, thick, trapezium, trapezium left angle=70,
               trapezium right angle=110, fill=black!8, align=center,
               minimum height=9mm, inner xsep=4mm, text width=40mm},
    op/.style={draw, thick, fill=akcent!10, align=center,
               minimum height=8mm, text width=54mm, inner sep=3pt},
    odrz/.style={draw, thick, fill=akcent!10, align=center,
                 minimum height=8mm, text width=31mm, inner sep=3pt},
    dec/.style={draw, thick, diamond, aspect=2.6, fill=white, align=center,
                inner sep=0pt, text width=42mm},
    strz/.style={->, thick},
    tn/.style={font=\scriptsize\bfseries, inner sep=1.5pt},
]
\hyphenpenalty=10000 \exhyphenpenalty=10000
\node[term] (start) at (0, 0) {START};
\node[we]   (wej)   at (0,-1.35) {Ramka DF17 z~pozycją, czas $t$};
\node[dec]  (d0)    at (0,-3.2) {Znana pozycja stacji?};
\node[op]   (opl)   at (0,-5.3)
     {$L\leftarrow$ pozycja lokalna z~jednej ramki (odniesienie: stacja)};
\node[op]   (opg)   at (0,-6.8)
     {$G\leftarrow$ pozycja globalna z~pary ramek (PipeDecoder) lub brak};
\node[op]   (tab)   at (0,-8.3)
     {Ustal $P$ i~$Q$ według tablicy decyzyjnej (tabela~\ref{tab:tablica_pozycja})};
\node[dec]  (d1)    at (0,-10.3) {Ustalono pozycję $P$?};
\node[dec]  (d2)    at (0,-13.1)
     {Stacja znana i~$P$ dalej niż \SI{450}{km} od stacji?};
\node[dec]  (d3)    at (0,-16.5)
     {Skok od poprzedniej pozycji $>\SI{600}{m/s}\cdot\Delta t+\SI{1}{km}$
      i~$Q\neq$ VERIFIED?};
\node[op]   (zap)   at (0,-19.0) {Zapisz $P$, $Q$ oraz czas pozycji $t$};
\node[term] (konz)  at (0,-20.3) {KONIEC};

\node[op, text width=26mm] (lbrak) at (6.2,-3.2) {$L\leftarrow$ brak};

\node[odrz] (r1) at (6.2,-10.3) {Odrzuć: brak pozycji lub niezgodne $L$ i~$G$};
\node[term] (k1) at (9.6,-10.3) {KONIEC};
\node[odrz] (r2) at (6.2,-13.1) {Odrzuć: poza zasięgiem};
\node[term] (k2) at (9.6,-13.1) {KONIEC};
\node[odrz] (r3) at (6.2,-16.5) {Odrzuć: nierealny skok pozycji};
\node[term] (k3) at (9.6,-16.5) {KONIEC};

\draw[strz] (start) -- (wej);
\draw[strz] (wej) -- (d0);
\draw[strz] (d0) -- node[tn, right=1mm] {TAK} (opl);
\draw[strz] (d0) -- node[tn, above, pos=0.35] {NIE} (lbrak);
\draw[strz] (opl) -- (opg);
\draw[strz] (lbrak.south) |- (opg.east);
\draw[strz] (opg) -- (tab);
\draw[strz] (tab) -- (d1);
\draw[strz] (d1) -- node[tn, right=1mm] {TAK} (d2);
\draw[strz] (d2) -- node[tn, right=1mm] {NIE} (d3);
\draw[strz] (d3) -- node[tn, right=1mm] {NIE} (zap);
\draw[strz] (zap) -- (konz);

\draw[strz] (d1) -- node[tn, above, pos=0.3] {NIE} (r1);
\draw[strz] (r1) -- (k1);
\draw[strz] (d2) -- node[tn, above, pos=0.3] {TAK} (r2);
\draw[strz] (r2) -- (k2);
\draw[strz] (d3) -- node[tn, above, pos=0.3] {TAK} (r3);
\draw[strz] (r3) -- (k3);
\end{tikzpicture}
\caption{Schemat blokowy wyznaczania pozycji statku powietrznego w~module
\texttt{adsb\_decoder.py}. Symbole: owal -- początek i~koniec algorytmu,
równoległobok -- dane wejściowe, prostokąt -- operacja, romb -- decyzja
z~wyjściami TAK i~NIE. Oznaczenia: $L$ -- pozycja wyznaczona metodą lokalną,
$G$ -- metodą globalną, $P$ -- pozycja wynikowa, $Q$ -- jakość pozycji
(LOCAL, GLOBAL lub VERIFIED), $\Delta t$ -- czas od poprzedniej pozycji tego
statku. Schemat obejmuje etap następujący po odfiltrowaniu ramek, które nie
są ramkami DF17 lub mają błędną sumę CRC}
\label{rys:schemat_pozycja}
\end{figure}

% ---------------------------------------------------------------------
\section{Dekodowanie ADS-B w~czasie rzeczywistym}
\label{sec:dekodowanie}
% Moduł adsb_decoder, węzeł decoder_node, dwie metody
% pozycji (lokalna z pozycją stacji, globalna z pary ramek), statusy
% LOCAL/GLOBAL/VERIFIED, kontrole wiarygodności, wyznaczanie odległości,
% azymutu i elewacji, publikacja AircraftState.

% ---------------------------------------------------------------------
\section{Rejestracja danych}
\label{sec:rejestracja}
% rosbag2 + MCAP, sesje nazywane czasem UTC, pliki co
% 10 min, odzyskiwanie po utracie zasilania (reindex).

% ---------------------------------------------------------------------
\section{Interfejs użytkownika}
\label{sec:interfejs}
% display_node (układ ekranu, kolory jakości pozycji,
% wskaźnik REC), skrypt przycisku zasilania.

% ---------------------------------------------------------------------
\section{Uruchamianie i~zamykanie systemu}
\label{sec:start_stop}
% start_nodes.sh, sygnały SIGTERM/SIGINT, trap, grupy
% procesów, SIGTTIN; napotkane problemy i ich rozwiązania.